% =====================================================================
%  THÈME 3 — Arrondir & opérations — CORRIGÉ — Niveau MOYEN
% =====================================================================
\documentclass[11pt,a4paper]{article}
\usepackage[utf8]{inputenc}
\usepackage[T1]{fontenc}
\usepackage{lmodern}
\usepackage[french]{babel}
\usepackage{amsmath,amssymb,mathtools}
\usepackage{icomma}
\usepackage{siunitx}
\sisetup{output-decimal-marker={,},per-mode=symbol}
\usepackage[version=4]{mhchem}
\usepackage{xcolor}
\usepackage[most]{tcolorbox}
\usepackage{enumitem}
\usepackage{fancyhdr}
\usepackage[a4paper,margin=2.2cm,headheight=26pt,headsep=8pt]{geometry}

\definecolor{primary}{RGB}{150,98,162}
\definecolor{primarydark}{RGB}{92,52,110}
\newcommand{\cs}{chiffre significatif}
\newcommand{\CS}{chiffres significatifs}

\pagestyle{fancy}\fancyhf{}
\fancyhead[L]{\small\textcolor{primarydark}{\textbf{Fiche 1} — Chiffres significatifs}}
\fancyhead[R]{\small\textcolor{primarydark}{\textsc{Corrigé · Moyen · Thème 3}}}
\fancyfoot[C]{\small\thepage}
\renewcommand{\headrulewidth}{0.8pt}
\renewcommand{\headrule}{\hbox to\headwidth{\color{primary}\leaders\hrule height \headrulewidth\hfill}}

\newtcolorbox[auto counter]{cor}[1][]{enhanced,breakable,boxrule=0.8pt,arc=2pt,
  colback=primarydark!5,colframe=primarydark,left=3mm,right=3mm,top=2mm,bottom=2mm,
  fonttitle=\bfseries,coltitle=white,
  title={Correction — Exercice~\thetcbcounter\ifx\relax#1\relax\relax\else\ ---\ \textmd{\itshape #1}\fi}}

\setlength{\parindent}{0pt}\setlength{\parskip}{3pt}

\begin{document}

\begin{center}
{\Large\bfseries\color{primarydark} Corrigé — Niveau Moyen}\\[1pt]
{\color{primary} Thème 3 : les deux règles des opérations}
\end{center}
\vspace{2mm}

\begin{cor}[Addition : la règle des décimales]
Valeur brute : $\num{12,4}+\num{7,25}+\num{103,682}=\num{123,332}$. La donnée la moins précise est
$\qty{12,4}{\milli\litre}$ ($1$ décimale) ; on arrondit donc à \textbf{1 décimale} :
\[ \boxed{V \approx \qty{123,3}{\milli\litre}}. \]
Inutile d'écrire les chiffres $3$ et $2$ finaux : $\num{12,4}$ ignore déjà tout au-delà du dixième.
\end{cor}

\begin{cor}[Soustraction : la perte de chiffres significatifs]
\begin{enumerate}[label=\alph*),leftmargin=1.6em,itemsep=1pt]
  \item $\num{12,47}-\num{12,3}=\num{0,17}$ ; la donnée à $1$ décimale ($\num{12,3}$) impose $1$
        décimale : $\boxed{\qty{0,2}{\gram}}$. On passe de $4$ et $3$~CS à \textbf{1}~CS.
  \item $\num{5,624}-\num{5,60}=\num{0,024}$ ; $\num{5,60}$ ($2$ décimales) impose $2$ décimales :
        $\boxed{\qty{0,02}{\gram}}$. On passe de $4$ et $3$~CS à \textbf{1}~CS.
\end{enumerate}
\textbf{Constat :} soustraire deux nombres \emph{voisins} détruit les CS. C'est le piège de la
« perte de chiffres significatifs », à fuir en chimie analytique.
\end{cor}

\begin{cor}[Multiplication et division : la règle des \CS]
\begin{itemize}[leftmargin=1.6em,itemsep=1pt]
  \item a) $\num{6,3}\times\num{2,1}=\num{13,23}$ ; min $=2$~CS $\Rightarrow \boxed{\num{13}}$.
  \item b) $\num{2,50}\times\num{3,142}=\num{7,855}$ ; min $=3$~CS ($\num{2,50}$) ; supprimé $5\,(\ge5)
        \Rightarrow \boxed{\num{7,86}}$.
  \item c) $\num{45,0}\div\num{1,2}=\num{37,5}$ ; min $=2$~CS ($\num{1,2}$) ; supprimé $5\,(\ge5)
        \Rightarrow \boxed{\num{38}}$.
\end{itemize}
\end{cor}

\begin{cor}[Quelle règle s'applique ?]
\begin{itemize}[leftmargin=1.6em,itemsep=1pt]
  \item a) $\num{2,5}+\num{13,42}$ : règle des \textbf{décimales} ; min $=1$ décimale $\to$ résultat à $1$ décimale.
  \item b) $\num{2,5}\times\num{13,42}$ : règle des \textbf{\CS} ; min $=2$~CS $\to$ résultat à $2$~CS.
  \item c) $\num{100,0}-\num{2,555}$ : règle des \textbf{décimales} ; min $=1$ décimale ($\num{100,0}$) $\to$ $1$ décimale.
  \item d) $\num{8,20}\div\num{4,1}$ : règle des \textbf{\CS} ; min $=2$~CS ($\num{4,1}$) $\to$ $2$~CS.
\end{itemize}
\end{cor}

\begin{cor}[Concentration massique]
\begin{enumerate}[label=\alph*),leftmargin=1.6em,itemsep=1pt]
  \item $c_m=\dfrac{\qty{4,80}{\gram}}{\qty{15,0}{\milli\litre}}=\num{0,32}\ \si{\gram\per\milli\litre}$.
        Les deux données ont $3$~CS ; on garde $3$~CS : $\boxed{c_m=\qty{0,320}{\gram\per\milli\litre}}$
        (le zéro de queue est obligatoire).
  \item $c_m=\qty{0,320}{\gram\per\milli\litre}\times\num{1000}=\qty{320}{\gram\per\litre}
        =\num{3,20e2}\ \si{\gram\per\litre}$ ($3$~CS).
\end{enumerate}
\end{cor}

\end{document}
